% ============================================================
% ch10.tex —— 第 10 章 Hitchin 表示
% 底本：LNM 2386《Spectral Networks》，印刷 pp.93–110（PDF p133–p150）
% ============================================================

% [原文 p.93]
\chapter{Hitchin 表示}

\ChapterAbstract{把 Fuchs 表示（Fuchsian representation）与 $\PSL_2(\R)$ 到伴随分裂实半单李群 $G$ 的 Kostant 主嵌入复合起来，给出一类具有深远重要性的基本群表示。所有这些表示都落在特征簇的同一个连通分支之中，即所谓的 Hitchin 分支，并且该分支中的所有表示都是离散且忠实的。特征簇的带标架版本允许一种簇结构，即具有特殊双有理转移映射的坐标系。这些转移映射尤其是不含减法的，由此给出「坐标为正的点」这一良定义的概念。特征簇的正部分正是 Hitchin 分支。}

\section{Hitchin 表示}

高阶泰希米勒空间的最初一批例子，早在 Fock 与 Goncharov 于 [129] 中引入这一术语之前就已经出现。这族例子涉及伴随分裂实半单李群 $G$ 的特征簇的连通分支，由 Hitchin 在 [199] 中引入。借助适当的嵌入 $\SL_2(\R)\hookrightarrow G$ 便可得到它们，而这一嵌入使得经典泰希米勒空间的许多基本拓扑与动力性质能够推广到这些高阶类似物。这些分支最初被定义为泰希米勒分支，但如今人们更常称它们为 Hitchin 分支，并称其中所含的基本群表示为 Hitchin 表示。本章将大量用到实李群与实李代数的理论。实李群的要点见 A.2 节。

% [原文 p.94]
\subsection{旗的构型}

现在我们把注意力转向分裂实李群 $\SL_n(\R)$，并描述旗簇的概念。对任意分裂实李群而言，这一概念都是刻画其 Hitchin 分支的核心；该概念的另一个例子将出现在 11.1.4 节。
考虑全旗簇，其定义为
\[
\mathcal{F}(\R^n) := \left\{F=(F_1,F_2,\dots,F_{n-1}) \mid F_i\subset\R^n,\ \dim(F_i)=i,\ F_i\subset F_{i+1}\right\}.
\]
若 $F_i\cap F'_{n-i}=\{0\}$，则称两个旗 $F,F'\in\mathcal{F}$ 是横截的。若对满足 $i+j+k=2n$ 的所有整数 $1\le i,j,k\le n$ 都有 $F_i\cap G_j\cap H_k=\{0\}$，则称旗三元组 $(F,G,H)$ 是横截的。特别地，这意味着对满足 $i+j+k=2n+1$ 的整数 $1\le i,j,k\le n$，空间 $F_i\cap G_j\cap H_k$ 是 $\R^n$ 中的一条直线。
对每个整数 $d\ge2$，用 $\mathrm{Conf}^d(\mathcal{F}(\R^n))$ 表示 $d$-元旗组的模空间，即 $(\mathcal{F}(\R^n))^d$ 用 $\SL_n(\R)$ 的对角作用取商。对称群 $\Sym_d$ 在 $(\mathcal{F}(\R^n))^d$ 上的自然作用下降到 $\mathrm{Conf}^d(\mathcal{F}(\R^n))$ 上。
我们将特别关心空间 $\mathrm{Conf}^2(\mathcal{F}(\R^n))$、$\mathrm{Conf}^3(\mathcal{F}(\R^n))$ 及其某些子空间。用 $\mathrm{Conf}^d_*(\mathcal{F}(\R^n))$ 表示由两两横截的 $d$-元旗组构成的构型空间。子空间 $\mathrm{Conf}^d_*(\mathcal{F}(\R^n))\subset\mathrm{Conf}^d(\mathcal{F}(\R^n))$ 在所有置换下不变。
一个简单的线性代数练习表明，群 $\SL_n(\R)$ 在 $\mathrm{Conf}^2_*(\mathcal{F}(\R^n))$ 上的作用是可迁的，并且 $\mathrm{Conf}^2_*(\mathcal{F}(\R^n))$ 中任一元素的稳定化子在 $\SL_n(\R)$ 中共轭于对角矩阵构成的子群。

% [Tutor Note: 全旗簇的定义只列到 $F_{n-1}$，按惯例约定 $F_0=\{0\}$、$F_n=\R^n$；横截条件 $F_i\cap F'_{n-i}=\{0\}$ 等价于 $F_i+F'_{n-i}=\R^n$，因为两者维数之和恰为 $n$。核对 $n=2$ 的情形：$\mathcal F(\R^2)$ 就是 $\R^2$ 中直线之集，两旗横截即对应两条不同直线，稳定化子为对角矩阵群，故 $\mathrm{Conf}^2_*$ 与 $\mathrm{Conf}^2$ 都只剩下一个点——这正是「$\SL_n(\R)$ 可迁作用」的最简单检验。]

\subsection{正旗}

现在固定 $\R^n$ 的一组标准基 $(e_1,\dots,e_n)$，令 $F\in\mathcal{F}$ 为满足 $F_i=\mathrm{Span}(e_1,\dots,e_i)$ 的旗，并令 $E\in\mathcal{F}$ 为满足 $E_i=\mathrm{Span}(e_n,\dots,e_{n-i+1})$ 的旗。于是任何与 $F$ 横截的旗 $T\in\mathcal{F}$ 都是 $E$ 在某个唯一幂幺矩阵 $u_T$ 下的像。一个幂幺矩阵称为全正的，当且仅当它的所有子式都为正，除了那些因矩阵幂幺这一条件而必须为零的子式之外。关于这类矩阵性质的更一般论述见 12.1 节。称旗三元组 $(E,T,F)$ 为正的，当且仅当矩阵 $u_T$ 是全正幂幺矩阵。注意任意两个横截的旗 $(F_1,F_2)$ 都可以被 $\SL_n(\R)$ 中某个元素映为旗对 $(E,F)$。因此上述正性概念可以推广为 $\mathcal{F}$ 中任意两两横截的旗三元组的正性。正旗三元组的模空间记作 $\mathrm{Conf}^3_+(\mathcal{F}(\R))$。由于任何正旗三元组自动是横截的，特别地其中各旗两两横截，故有 $\mathrm{Conf}^3_+(\mathcal{F}(\R))\subset\mathrm{Conf}^3_*(\mathcal{F}(\R))$。若存在全正幂幺矩阵 $u_1,\dots,u_{k-2}$，使得旗 $k$-元组 $(F_1,\dots,F_k)$ 可以被映到 $(F,u_1F,u_1u_2F,\dots,u_1\cdots u_{k-2}F,E)$，则称该旗 $k$-元组是正的。

% [Tutor Note: 「除幂幺条件迫使为零的子式之外」这一限定不可省：幂幺矩阵 $u_T$ 的对角元全为 1，而下三角位置（$i>j$）必为零。若对这些位置也要求「子式为正」就会自相矛盾。以 $n=3$ 为例，$1\times1$ 子式即矩阵元本身，其中三个下三角元为零、不参与正性要求；$2\times2$ 子式共 9 个，凡取到两个下三角元或取整个下三角块的均为零（例如 $\det\begin{pmatrix}u_{21}&u_{22}\\u_{31}&u_{32}\end{pmatrix}=-u_{32}$ 之类要逐个核对），只需检验剩余那些其值由幂幺性决定非零的子式。]

% [原文 p.95]
旗簇 $\mathcal{F}_G$ 中旗组的正性可以对任意分裂实李群 $G$、按与 $\SL_n(\R)$ 情形类似的方式借助 $G$ 的 Borel 子群来定义，这是 Lusztig [251] 的工作。

\subsection{到伴随分裂实李群的 Hitchin 表示}

设 $S_g$ 为闭定向曲面。对于伴随分裂实半单李群 $G$，存在唯一嵌入 $\pi:\PSL_2(\R)\to G$，它是与 $\g$ 的一个主三维子代数相对应的李群同态，该子代数称为 Kostant 主子代数 $\sl_2(\R)\subset\g$，由 [239, Lem. 5.2] 引入（完整描述见定义 A.4）。
对一个固定的离散嵌入 $\iota:\pi_1(S_g)\to\PSL_2(\R)$，Hitchin 在 [199] 中证明：包含主 Fuchs 表示 $\pi\circ\iota:\pi_1(S_g)\to G$ 的子空间是一个连通分支，而且事实上是拓扑平凡的，其维数为 $(2g-2)\dim_{\R}\g$。在群为 $G=\PSL_2(\R)$ 的特殊情形，该分支就是泰希米勒空间。具体地，我们有：

\begin{theorem}[10.1] [199, Thm. A] 设 $S_g$ 是亏格 $g\ge2$ 的紧致闭连通拓扑曲面，$G$ 是伴随分裂实半单李群。则特征簇 $\mathcal{R}(S_g,G)$ 有一个连通分支同胚于维数为 $(2g-2)\dim_{\R}\g$ 的欧几里得球。
\end{theorem}

于是我们如下定义 Hitchin 分支：

\begin{definition}[10.2] 设 $G$ 为伴随分裂实半单李群，$S_g$ 为亏格 $g\ge2$ 的紧致闭连通拓扑曲面。则 Hitchin 分支 $T_H(S_g,G)$ 是特征簇 $\mathcal{R}(S_g,G)$ 中包含主 Fuchs 表示 $\pi\circ\iota:\pi_1(S_g)\to G$ 的连通分支。
\end{definition}

Hitchin 分支中所含的基本群表示在 [129, 242] 中被证明是离散且忠实的。而在 $\PSL_n(\R)$ 的特殊情形下，由上面讨论的正旗还得到了 Hitchin 分支的一个非常有用的刻画：

\begin{theorem}[10.3] [129, 175, 242] 设 $\rho:\pi_1(S_g)\to G$ 是一个表示，$\mathcal{F}_G$ 表示任意分裂实李群 $G$ 的旗空间。则 $\rho\in T_H(S_g,G)$ 当且仅当存在连续的 $\rho$-等变映射 $\xi:\partial\pi_1(S_g)\to\mathcal{F}_G$，它把 $\partial\pi_1(S_g)$ 中的正元组映为 $\mathcal{F}_G$ 中的正元组。
\end{theorem}

\begin{remark}[10.4] 若 $\bar S_{g,h}$ 是亏格为 $g$、带有 $n$ 个边界分支的紧致曲面，则可以给出类似的构造。我们需要固定一个离散嵌入 $\iota:\pi_1(\bar S_{g,h})\to\PSL_2(\R)$，它来自 $\bar S_{g,h}$ 上具有测地边界的双曲结构。然而此时若 $G$ 连通，则 $\mathcal{R}(\bar S_{g,h},G)$ 是连通空间，我们无法直接应用上面的构造。在这种情形下，可以考虑 $\mathcal{R}(\bar S_{g,h},G)$ 的如下子空间：设 $\mathcal{R}^p(\bar S_{g,h},G)$ 为所有这样的表示
% [原文 p.96]
$\rho$ 构成的空间：对每个同伦于某一边界分支的 $\gamma\in\pi_1(\bar S_{g,h})$，$\rho(\gamma)$ 固定 $\mathcal{F}_G$ 中的一个元素。这类表示称为外周抛物的（peripherally parabolic）。然后令 $T_H(\bar S_{g,h},G)$ 为 $\mathcal{R}^p(\bar S_{g,h},G)$ 中包含 $\pi\circ\iota:\pi_1(\bar S_{g,h})\to G$ 的连通分支。则 $T_H(\bar S_{g,h},G)$ 由离散且忠实的表示构成 [131]。
\end{remark}

% [Tutor Note: 注 10.4 的记号与文字不一致：原文记号写作 $\bar S_{g,h}$（下标 $h$ 按惯例表示边界分支数），正文却说「$n$ 个边界分支」。此处应理解为「$h$ 个边界分支」，两种写法不应同时出现；下文同样用 $h$ 作下标。这是原文的记号笔误。]

\begin{remark}[10.5] 映射 $\xi$ 把 $\partial\pi_1(S_g)$ 中的正元组映为 $\mathcal{F}_G$ 中的正元组，当且仅当它把 $\partial\pi_1(S_g)$ 中的正四元组映为 $\mathcal{F}_G$ 中的正四元组。
\end{remark}

\section{Fock–Goncharov 坐标}

\subsection{横截与正带标架表示}

模空间理论中有一条一般原则：若给所关心的结构再添加一些附加数据，则该结构的模空间可能反而更容易描述。这一原则在本节中得到了体现：我们将用一份称为标架的附加数据来丰富我们的表示。虽然带标架表示的模空间看起来可能比特征簇更复杂，但利用 Fock 与 Goncharov 在 [129] 中引入的特殊坐标系，前者反而更容易描述。自此以后，关于原特征簇的信息都可以通过遗忘标架而得到。
设 $S$ 是非紧纤毛曲面，赋予它具有测地边界的完备有限体积双曲结构。设 $\tilde S\subseteq\HH^2$ 是 $S$ 的万有覆叠，并设 $\tilde P\subseteq\partial_\infty^{\,c,d}\tilde S$ 是 $S$ 的所有穿孔与纤毛在 $\tilde S$ 边界上的提升之并。设 $G$ 是带有 Borel 子群 $B$ 的分裂实李群，并设 $\mathcal{F}_G=G/B$ 为全旗簇。

\begin{definition}[10.6] 一个（旗）标架是一个映射 $F:\tilde P\to\mathcal{F}_G$。设 $\rho:\pi_1(S)\to G$ 是同态。$\rho$ 的一个标架是指一个 $\pi_1(S)$-等变的标架 $F:\tilde P\to\mathcal{F}_G$，即对每个 $\gamma\in\pi_1(S)$ 与所有 $\tilde p\in\tilde P$ 有 $F(\gamma\tilde p)=\rho(\gamma)F(\tilde p)$。带标架同态是指一对 $(\rho,F)$，其中 $F$ 是 $\rho$ 的标架。
所有带标架同态构成的空间记作 $\Hom^{\mathrm{fr}}(\pi_1(S),G)$。空间
\[
\mathcal{E}^{\mathrm{fr}}(S,G) := \Hom^{\mathrm{fr}}(\pi_1(S),G)/G
\]
称为带标架表示的模空间。带标架表示是 $\mathcal{E}^{\mathrm{fr}}(S,G)$ 中的一个元素。带标架特征簇是带标架可约表示空间的下述豪斯多夫商
\[
\mathcal{R}^{\mathrm{fr}}(S,G) := \Hom^{\mathrm{red},\mathrm{fr}}(\pi_1(S),G)/G.
\]
\end{definition}

\begin{remark}[10.7] 注意当 $G$ 是实李群时，并非每个同态 $\rho:\pi_1(S)\to G$ 都允许标架。允许标架的同态称为外周抛物的。
\end{remark}

% [原文 p.97]
设 $\Delta$ 是 $S$ 的一个理想三角剖分，$\tilde\Delta$ 表示 $\Delta$ 到万有覆叠 $\tilde S$ 的提升。

\begin{definition}[10.8] 带标架同态 $(\rho,F)$ 称为 $\Delta$-横截的，若对 $\tilde\Delta$ 的每个以 $\tilde p_1,\tilde p_2,\tilde p_3\in\tilde P$ 为顶点的三角形，旗三元组 $(F(\tilde p_1),F(\tilde p_2),F(\tilde p_3))$ 都是横截的。
所有 $\Delta$-横截带标架同态构成的空间记作 $\Hom^{\mathrm{fr}}_\Delta(\pi_1(S),G)$。空间
\[
\mathcal{E}^{\mathrm{fr}}_\Delta(S,G) := \Hom^{\mathrm{fr}}_\Delta(\pi_1(S),G)/G
\]
称为 $\Delta$-横截带标架表示的模空间。
\end{definition}

\begin{definition}[10.9] 标架 $F$ 称为正的，若对每个循环定向的元组 $(\tilde p_1,\dots,\tilde p_k)\in\tilde P^k$，旗元组 $(F(\tilde p_1),\dots,F(\tilde p_k))$ 都是正的（按 10.1.2 节的意义）。同态 $\rho\in\Hom(\pi_1(S),G)$ 称为正的，若它允许一个正标架 $F$。称对 $(\rho,F)$ 为正带标架同态。
所有正带标架同态构成的空间记作 $\Hom^{\mathrm{fr}}_+(\pi_1(S),G)$。空间
\[
\mathcal{R}^{\mathrm{fr}}_+(S,G) := \Hom^{\mathrm{fr}}_+(\pi_1(S),G)/G
\]
称为正带标架表示的模空间。
\end{definition}

\begin{remark}[10.10] 由该定义立即推出，正带标架表示关于 $S$ 的任一理想三角剖分都是横截的（参见 10.1.2 节）。
\end{remark}

\begin{remark}[10.11] 注意，一般而言 $G$ 在 $\Hom(\pi_1(S),G)$ 上与在 $\Hom^{\mathrm{fr}}(\pi_1(S),G)$ 上的作用并非所有轨道都闭。因此商空间 $\mathcal{E}(S,G)$ 与 $\mathcal{E}^{\mathrm{fr}}(S,G)$ 不是豪斯多夫的。然而 $\mathcal{E}^{\mathrm{fr}}(S,G)$ 中正带标架表示的空间 $\mathcal{R}^{\mathrm{fr}}_+(S,G)$ 构成一个豪斯多夫子空间。
此外，自然的遗忘映射 $\mathcal{E}^{\mathrm{fr}}(S,G)\to\mathcal{E}(S,G)$ 把正带标架表示的空间 $\mathcal{R}^{\mathrm{fr}}_+(S,G)$ 满射到 10.1 节讨论的 Hitchin 分支 $T_H(S,G)$ 上，且纤维有限。这意味着每个 Hitchin 表示都允许标架；事实上它至多允许有限多个标架。注意在一般情形下这是不对的：存在不允许标架的表示，也存在允许无穷多个标架的表示（更多细节参见 [131]）。
\end{remark}

事实上，为检验一个标架是否为正，只需对三角剖分的所有正定向四边形检验，而不必对 $\tilde P$ 中所有可能的循环定向元组检验：

\begin{proposition}[10.12] ([131]) 设 $\Delta$ 是 $S$ 的三角剖分，$F$ 是标架，且使得对 $\tilde\Delta$ 的每个正定向四边形 $(\tilde p_1,\tilde p_2,\tilde p_3,\tilde p_4)$，旗四元组 $(F(\tilde p_1),F(\tilde p_2),F(\tilde p_3),F(\tilde p_4))$ 都是正的。则标架 $F$ 是正的。
\end{proposition}

\begin{remark}[10.13] 关于正标架更一般的讨论，参见 [178]。
\end{remark}

% [原文 p.98]
\begin{definition}[10.14] $S$ 的 $G$-高阶带标架泰希米勒空间（或泰希米勒 $X$-空间）是定义 10.9 中定义的正带标架表示的模空间 $\mathcal{R}^{\mathrm{fr}}_+(S,G)$。
\end{definition}

称这样的空间为高阶泰希米勒空间，其理由由下述命题给出。

\begin{proposition}[10.15]（[129] 中的定理 1.9 与 1.10）任何正表示 $\rho\in\mathcal{R}^{\mathrm{fr}}_+(G,S)$ 都是离散且忠实的。
\end{proposition}

$\mathcal{R}^{\mathrm{fr}}_+(S,G)$ 在注 10.11 的自然遗忘映射下的像，是 $\pi_1(S)$ 到 $G$ 的一族只由离散且忠实表示构成的表示空间，即它也是一个高阶泰希米勒空间，只不过它只描述基本群的表示。不过，考察 $\mathcal{R}^{\mathrm{fr}}_+(S,G)$ 而非后一个空间是有道理的，因为 $\mathcal{R}^{\mathrm{fr}}_+(S,G)$ 允许自然的坐标，这些坐标既容易定义，本身也很有趣。

\begin{remark}[10.16] 当 $G=\PSL_2(\R)$ 时，$\mathcal{R}^{\mathrm{fr}}_+(S,\PSL_2(\R))$ 在遗忘映射下的像正是最自然地称为 $S$ 的泰希米勒空间的那个空间。然而当 $S$ 非闭时，还存在 $S$ 的泰希米勒空间的其他（略有不同的）版本，例如 $\mathcal{R}^{\mathrm{fr}}_+(S,\PSL_2(\R))$。后者是定义 7.3 中引入的带标架泰希米勒空间 $T^x(S)$，并自然地由剪切坐标参数化。当 $G=\SL_2(\R)$ 时，泰希米勒空间的另一个例子是由 Penner [297, 298] 引入的所谓装饰泰希米勒空间，它带有自然坐标，这些坐标赋予该空间 A-型簇代数簇的结构（参见 5.5 节）。本书只考虑带标架泰希米勒空间的高阶推广，即空间 $\mathcal{R}^{\mathrm{fr}}_+(S,G)$，而装饰泰希米勒空间的高阶推广的定义可参见 [129]。
\end{remark}

\subsection{带标架表示空间的参数化}

我们要研究的 $\mathcal{R}^{\mathrm{fr}}(S,G)$（或更确切地说它的一个开稠密子集）上自然坐标之构造的基石，是下述分解定理。它表明：当 $G$ 的中心平凡时，给定 $S$ 的三角剖分 $\Delta$，在 $\mathcal{R}^{\mathrm{fr}}_+(S,G)$ 的一个开稠密子集上，可以把一个正带标架表示完全描述为它在 $\Delta$ 的每个面——三角形 $T$——上的限制数据，再加上粘合数据：在 $\Delta$ 的每条两侧面被粘合的边上，给出 $G$ 的 Cartan 子群中的一个元素。

\begin{theorem}[10.17]（分解定理（[129]，第 6 节））设 $\Delta$ 是 $S$ 的理想三角剖分，$F(\Delta)$（相应地 $E(\Delta)$）是 $\Delta$ 的面（相应地边）的集合。若 $G$ 的中心平凡，则带标架表示的模空间 $\mathcal{R}^{\mathrm{fr}}(S,G)$ 双有理等价于
\begin{equation}
\prod_{T\in F(\Delta)}\mathcal{R}^{\mathrm{fr}}(T,G)\;\times\;\prod_{e\in E(\Delta)}H, \tag{10.1}
\end{equation}
其中 $H$ 是 Borel 子群 $B$ 对应的 Cartan 子群。
\end{theorem}

% [原文 p.99]
\begin{remark}[10.18] 有理映射与双有理等价的概念本质上是代数性的。粗略地说，两个空间双有理等价意味着它们同构到相差正余维子簇为止。这里所说的子簇由非一般的带标架表示构成，例如两个不同穿孔 $\tilde P$ 在 $\mathcal{F}_G$ 中的像重合的那些表示。关于代数几何与有理映射的更多细节可参见 [189]。
\end{remark}

模空间 $\mathcal{R}^{\mathrm{fr}}(T,G)=\mathcal{E}^{\mathrm{fr}}(T,G)$（其中 $T$ 是理想三角形）可与旗三元组的构型空间等同。因此我们需要在旗三元组空间上定义坐标，以及对应于三角剖分之边的坐标。我们不会完整地描述这一构造，读者可参阅 [129]。我们将在 10.2.3 节详细给出当 $G=\PGL_n(\R)$ 时坐标的构造，并在 10.2.4 节说明如何由这些坐标恢复带标架表示。这里我们专注于给出 $\mathcal{R}^{\mathrm{fr}}(S,G)$ 上坐标的一般性质。这些坐标赋予带标架表示空间一个 X-型簇代数簇的结构（见第 5 章）：

\begin{theorem}[10.19] ([131]) 对任意纤毛曲面 $S$ 与任意无中心的分裂实李群 $G$，存在一个种子 $\mathcal{S}$（按定义 5.2 的意义）以及一个映射
\begin{equation}
\mathcal{R}^{\mathrm{fr}}(S,G)\to\mathcal{X}_{\mathcal{S}}, \tag{10.2}
\end{equation}
它是双有理同构。
\end{theorem}

种子 $\mathcal{S}$ 依赖于 $G$ 以及 $S$ 的理想三角剖分 $\Delta$ 的选取，但由 $\Delta$ 经过一次翻转得到的三角剖分所给出的种子，可由 $\mathcal{S}$ 出发经一串变异得到，因此给出同一个簇代数簇。于是我们可以把这个簇代数簇记作 $\mathcal{X}_{\mathcal{S}}=\mathcal{X}_{S,G}$。注意在多数情形下，经过一串变异后得到的种子并不对应于 $S$ 的某个三角剖分。由于 $\mathcal{R}^{\mathrm{fr}}(S,G)$ 是 X-型簇代数簇，它带有典范的泊松结构。

回忆第 5 章：在实簇代数簇中，正点的轨迹是良定义的。在特征簇中，这一轨迹恰好就是正带标架表示的空间：

\begin{theorem}[10.20] ([131]) 空间 $\mathcal{R}^{\mathrm{fr}}_+(S,G)$ 微分同胚于 $\mathcal{X}_{S,G}$ 的实正点集。
\end{theorem}

随后，对坐标取对数便知，$G$-高阶泰希米勒空间 $\mathcal{X}^+_{S,G}$ 同胚于开球 $\R^N$，其中 $N\in\R$。

% [Tutor Note: 原文此处写「$N\in\R$」应为笔误，正确形式是 $N\in\N$：$\R^N$ 中的 $N$ 是空间的维数，由分解定理 (10.1) 它等于 $\#F(\Delta)\cdot\dim\mathcal{R}^{\mathrm{fr}}(T,G)+\#E(\Delta)\cdot\dim H$，显然是有限非负整数。判断依据：同胚于 $\R^N$ 的说法要求 $N$ 为非负整数，$N\in\R$ 会使 $\R^N$ 无意义。]

\begin{remark}[10.21] 对单连通的李群 $G$，有一个类似的构造，给出模空间 $\mathcal{A}_{S,G}$，它带有 A-型簇代数簇的结构。它称为扭装饰表示的模空间。
\end{remark}

% [原文 p.100]
有一个重要的技术性复杂之处：正如我们在由 $n$ 个向量构成的模空间这一基本例 5.14 中所看到的，存在一个扭转移位映射 $\tau:(v_1,v_2,\dots,v_n)\mapsto(v_2,\dots,v_n,s_G v_1)$，其中 $s_G$ 是 $G$ 中一个非常特殊的二阶元。对 $G=\SL_m(\C)$，当 $m$ 为偶数时它是 $-\mathrm{id}$，当 $m$ 为奇数时它是 $\mathrm{id}$。
为在局部系统中实现该扭转移位，我们考虑带孔切丛 $T'S:=TS\setminus S$。它的基本群进入短正合列
\[
1\to\Z\to\pi_1(T'S)\to\pi_1(S)\to1.
\]
$S$ 上的扭曲 $G$-局部系统是指 $T'S$ 上的 $G$-局部系统，其在任一纤维中绕生成元的单值为 $s_G$。最后，我们给扭曲 $G$-局部系统 $L$ 装饰一个装饰旗的选取，即在纤毛曲面 $S$ 的每个边界分支上取一个完全旗，并在每个子空间上给定一个（复）体积形式。装饰旗可以看作一个元组 $(v_1,v_1\wedge v_2,v_1\wedge v_2\wedge v_3,\dots,\omega)$，其中 $\omega$ 是一个固定体积形式（它给出结构群到 $\SL_m(\C)$ 的约化）。更一般地，它是陪集空间 $G/U$ 中的一个元素，其中 $U$ 是极大幂幺子群。
对 $G=\SL_m(\C)$，装饰扭 $G$-局部系统的模空间 $\mathcal{A}_{S,G}$ 上的簇结构按与 X-坐标类似的方式得到：从理想三角剖分出发并考虑它的 $m$-细分。在一个三角形内，其外侧顶点上分别有装饰旗 $(f_1,f_1\wedge f_2,\dots)$、$(g_1,g_1\wedge g_2,\dots)$ 与 $(h_1,h_1\wedge h_2,\dots)$，我们对每个重心坐标为 $(a,b,c)$（其中 $a+b+c=m$）的内顶点以
\[
\Delta_{(a,b,c)} := \omega(f_1\wedge f_2\wedge\dots\wedge f_a\wedge g_1\wedge\dots\wedge g_b\wedge h_1\wedge\dots\wedge h_c)
\]
赋予一个坐标。所有 $\Delta_{(a,b,c)}$ 构成 $\mathcal{A}_{S,\SL_m(\C)}$ 的一个正正则图册。更多细节参见 [129, 第 8 章]。

\begin{remark}[10.22] 模空间 $\mathcal{X}_{S,G}$ 与 $\mathcal{A}_{S,G}$ 构成与同一变异类种子相关联的一对簇代数簇；这样一对空间称为簇系综（cluster ensemble）。在此语境下，5.5 节讨论的簇对偶性断言：整热带点 $\mathcal{A}_{S,G}(\Z^t)$ 参数化 $\mathcal{X}_{S,G}$ 上一族重要函数的一组基，反之亦然。
\end{remark}

\subsection{当 $G=\mathrm{PGL}_n$ 时的特殊坐标}

本部分中，$\mathbb{K}$ 一律表示 $\R$ 或 $\C$。
本部分将详细给出 $\mathcal{R}(S,\PGL_n(\mathbb{K}))$ 上簇坐标的构造。当 $\mathbb{K}=\R$ 时，这推广了第 7 章给出的 $T^x(S)$ 上坐标的构造。这一构造将是第 5 章所定义的 X-型簇代数簇的典型例子。设 $\Delta$ 是 4.3 节所定义的纤毛曲面 $S$ 的三角剖分，且 $\Delta$ 不含任何自折三角形，也不含任何二边形（见图 10.1）。

% [原文 p.101]
\Fig[0.72\linewidth]{figures/fig_10_1.png}{10.1}{左边是一个自折三角形，右边是一个二边形。}

设 $\mathrm{T}$ 是 $\R^3$ 中的欧几里得三角形，即坐标为 $A=(1,n,n)$、$B=(n,1,n)$ 与 $C=(n,n,1)$ 的三点的凸包。$\mathrm{T}$ 中形如 $(a,b,n)$ 且 $a+b=n+1$ 的点集（相应地，形如 $(n,b,c)$ 且 $b+c=n+1$ 的点集，形如 $(a,n,c)$ 且 $a+c=n+1$ 的点集）称为边 $AB$（相应地 $BC$、$AC$），或称为与 $C$（相应地 $B$、$A$）相对的边。这个欧几里得三角形 $\mathrm{T}$ 可以像图 10.2 那样分解为更小的三角形。

\Fig[0.85\linewidth]{figures/fig_10_2.png}{10.2}{从左到右：$\mathrm{T}$ 的一个 4-细分、一个 5-细分与一个 6-细分。}

这些更小三角形的顶点是坐标为 $(a,b,c)$、$a+b+c=2n+1$ 的点，其中 $a,b,c\ge1$ 为整数；每个小三角形属于两种类型之一：朝上的三角形以 $(a-1,b,c)$、$(a,b-1,c)$、$(a,b,c-1)$ 为顶点，其中 $a+b+c=2n+2$；朝下的三角形以 $(a+1,b,c)$、$(a,b+1,c)$、$(a,b,c+1)$ 为顶点，其中 $a+b+c=2n$。我们称这种把 $\mathrm{T}$ 分成小三角形的分法为 $\mathrm{T}$ 的 $n$-细分。在 $\mathrm{T}$ 的每条边上都有 $n-1$ 个更小的朝上三角形（在图 10.2 中以灰色标出）。对 $\Delta$ 的每个三角形 $T$，我们固定一个从 $\mathrm{T}$ 到 $T$ 的微分同胚，使得对 $\Delta$ 的所有边 $\gamma$，与 $\gamma$ 相邻的两个三角形的 $n$-细分在 $\gamma$ 上一致。
现在我们如下定义嵌入 $S$ 中的箭图 $Q_{n,\Delta}$：$Q_{n,\Delta}$ 的顶点是 $\Delta$ 的每个三角形的 $(n+1)$-细分的所有顶点，但除去那些本身是 $S$ 的穿孔的顶点（即 $\Delta$ 的三角形本身的顶点）。对 $(n+1)$-细分的每个朝下三角形，在每两个顶点之间有一支重数为 1 的箭，使得小三角形的边界按顺时针方向定向，见图 10.3。$Q_{n,\Delta}$ 中位于 $\Delta$ 的某条同伦于 $S$ 的边界分支的边上的顶点是冻结的。
我们在第 5 章已经看到，这个箭图给出一个种子 $\mathcal{S}_\Delta$，它定义了一个带有坐标函数的 X-环面 $\mathcal{X}_\Delta$。下面我们将说明，一个 $\Delta$-横截的带标架表示 $(\rho,F)$ 对应于 $\mathcal{X}_\Delta$ 的一个点。为此我们将构造对应于 $(\rho,F)$ 的 X-坐标。我们用 $\mathcal{X}_{S,\PGL_n(\mathbb{K})}$ 表示与初始种子 $\mathcal{S}_\Delta$ 相关联的 X-簇代数簇。一个带标架表示
% [原文 p.102]
\Fig[0.85\linewidth]{figures/fig_10_3.png}{10.3}{从左到右：$n=2,3,4$ 时，$\Delta$ 的一个三角形内的箭图 $Q_{n,\Delta}$。}
对应于 $\mathcal{X}_{S,G}$ 中的一个点，仅当它对 $S$ 的某个三角剖分 $\Delta$ 是 $\Delta$-横截的。
设 $(\rho,F)$ 是 $\Delta$-横截的带标架表示。我们从与位于 $\Delta$ 的三角形 $T$（顶点为 $p_1,p_2,p_3$）内部的 $Q_{n,\Delta}$ 的顶点相对应的那些坐标开始。我们首先为 $\mathbb{K}^3$ 中的旗三元组定义一个称为三重比的概念。设 $(A,B,C)$ 是 $\mathbb{K}^3$ 中的横截旗三元组，$\omega$ 是 $\mathbb{K}^3$ 上的体积形式。设 $a_1\in A_1$、$a_2\in A_2/A_1$、$b_1\in B_1$、$b_2\in B_2/B_1$、$c_1\in C_1$、$c_2\in C_2/C_1$ 均非零。旗三元组 $(A,B,C)$ 的三重比定义为
\[
r(A,B,C) = \frac{\omega(a_1,a_2,b_1)\,\omega(b_1,b_2,c_1)\,\omega(c_1,c_2,a_1)}{\omega(a_1,a_2,c_1)\,\omega(b_1,b_2,a_1)\,\omega(c_1,c_2,b_1)}.
\]
注意三重比不依赖于 $a_1,a_2,b_1,b_2,c_1,c_2$ 的选取，也不依赖于体积形式 $\omega$ 的选取。

% [Tutor Note: 三重比的这两条「无关性」可以逐项核对：$a_2$ 只在 $A_2/A_1$ 中有定义，若把它换成 $a_2+\alpha a_1$，则分子分母中含 $a_2$ 的因子各变为 $\omega(a_1,a_2,b_1)+\alpha\omega(a_1,a_1,b_1)$、$\omega(a_1,a_2,c_1)+\alpha\omega(a_1,a_1,c_1)$，而后两项为 $0$，故比值不变；其余 $b_2,c_2$ 同理。再把 $\omega$ 换成 $k\omega$，分子与分母各为三个括号之积，同时乘以 $k^3$，比值也不变。]

设 $(A,B,C)$ 是 $\mathbb{K}^n$ 中的横截旗三元组，$1\le a,b,c\le n$ 是满足 $a+b+c=2n+3$ 的整数。空间 $P_{a,b,c}=A_a\cap B_b\cap C_c$ 的维数为 3，并且旗 $A$、$B$、$C$ 在 $P_{a,b,c}$ 上诱导出一个横截旗三元组。我们用 $r_{a,b,c}(A,B,C)$ 表示相应的交比。由于 $Q_{n,\Delta}$ 在 $\Delta$ 的三角形 $T$ 中的非冻结顶点以 $(n+1)$-细分中坐标为 $(a,b,c)$、$a+b+c=2n+3$ 的点来标号，我们可以对这样的顶点赋予三重比 $r_{a,b,c}(F(p_1),F(p_2),F(p_3))$。
现在考虑 $Q_{n,\Delta}$ 中位于 $\Delta$ 的内部边 $\gamma$ 上的顶点，设 $p_1$ 与 $p_2$ 是 $\gamma$ 的端点。设 $p_3$ 与 $p_4$ 是与围绕 $\gamma$ 的两个三角形 $T_1$ 与 $T_2$ 相关联的其余顶点，如图 10.4。$Q_{n,\Delta}$ 中位于 $\gamma$ 上的顶点 $v$ 在三角形 $T_1$（相应地 $T_2$）中的坐标为 $(a,b,n+1)$（相应地 $(a,n+1,b)$），其中 $a+b=n+2$。空间 $P_{a,b}=F_a(p_1)\cap F_b(p_2)$ 是一个平面，四个旗 $F(p_1),F(p_2),F(p_3),F(p_4)$ 在 $P_{a,b}$ 中诱导出四条直线。用 $r_{a,b}(F(p_1),F(p_2),F(p_3),F(p_4))$ 表示与这个四元组相关联的交比，并把它赋给顶点 $v$。$\Delta$ 的外边上的坐标是冻结的，取为 1。
可以证明三重比与交比遵循第 5 章给出的 X-变异公式，而我们将于 10.2.4 节说明可以由 X-坐标重构带标架表示，从而得到下述定理：

% [原文 p.103]
\Fig[0.45\linewidth]{figures/fig_10_4.png}{10.4}{三角剖分的一条边 $\gamma$ 包含于两个三角形 $T_1$ 与 $T_2$ 之中，周围有穿孔 $p_1,p_2,p_3,p_4$。}

\begin{theorem}[10.23] ([131]) 映射
\begin{equation}
\mathcal{R}^{\mathrm{fr}}(S,\PGL_n(\mathbb{K}))\to\mathcal{X}_{S,\PGL_n(\mathbb{K})}, \tag{10.3}
\end{equation}
把带标架表示映到上述构造的坐标，是一个双有理同构。
\end{theorem}

不过要注意，在多数情形下，经过一串变异后得到的并不是 $S$ 的某个三角剖分的细分，从而这些坐标不能表示为 $S$ 的穿孔与纤毛处旗的三重比与交比。然而这些坐标卡确实存在：它们由从 $S$ 的三角剖分之细分出发的一系列变异所定义。

\subsection{蛇}

本节中 $\mathbb{K}$ 一律表示 $\R$ 或 $\C$。
蛇是由 Fock 与 Goncharov 在 [131] 中引入的组合工具，用来描述与一个旗三元组相伴的射影基，以便由 X-坐标的数据重构曲面 $S$ 的基本群的带标架 $\PGL_n(\mathbb{K})$-表示。存在蛇的变换这一概念，它诱导出相伴射影基的易于计算的变化。
让我们先引入蛇及其变换。设 $n\in\N$。回忆 10.2.3 节中给出的欧几里得三角形 $\mathrm{T}$ 的 $n$-细分的定义。我们把该 $n$-细分看作 $\mathrm{T}$ 上的一个（无向）图 $S_n$。

\begin{definition}[10.24] $S_n$ 上以 $A$（相应地 $B$、$C$）为基点的蛇是一条长为 $n$ 的定向道路 $A=v_1,\dots,v_n$，它从 $A$（相应地 $B$、$C$）出发，终止于边 $BC$（相应地 $AC$、$AB$）上的一个顶点 $v_n$（图 10.5）。
\end{definition}

\begin{remark}[10.25] 特别地，以 $A$ 为基点的蛇的各顶点之第一坐标构成的序列是递增的。
\end{remark}

% [Tutor Note: 注 10.25 可直接核对：以 $A=(1,n,n)$ 为基点时，沿蛇每走一步第一坐标增加 1，走 $n-1$ 步后到达第一坐标为 $n$ 的格点，而边 $BC$ 恰好是「第一坐标等于 $n$」的那些点（见 10.2.3 节对边的描述），故第一坐标序列为 $1,2,\dots,n$，严格递增。]

设 $(P,Q)\in\{A,B,C\}^2$，$P\ne Q$，并设 $s_{PQ}$ 表示以 $P$ 为基点、终止于 $Q$ 的唯一一条蛇。这些蛇都停留在边 $PQ$ 之内。

% [原文 p.104]
\Fig[0.55\linewidth]{figures/fig_10_5.png}{10.5}{以 $A$ 为基点的 $S_5$ 上的一条蛇。}

设 $(F^A,F^B,F^C)\in\mathrm{Conf}^3(\mathcal{F}(\mathbb{K}^n))$ 是 $\mathbb{K}^n$ 中的横截旗三元组。回忆 $\PGL_n(\mathbb{K})$ 作用在 $\mathbb{K}^n$ 的旗三元组上。所谓射影基，是指 $\mathbb{K}^n$ 的一个线性基模去把所有向量同乘一个标量。射影基与 $\PGL_n(\mathbb{K})$ 的元素一一对应。
对 $S_n$ 中坐标为 $(a,b,c)$ 的顶点，我们赋予空间 $F_{a,b,c}=F^A_a\cap F^B_b\cap F^C_c$；由于旗三元组横截且 $a+b+c=2n+1$，它是一条直线。
设 $e$ 是 $S_n$ 的一条有向边。恰有一个朝上的三角形 $t$ 与 $e$ 相邻。设 $V_1$（相应地 $V_2$）是与 $e$ 的起点（相应地终点）相关联的直线，$V_3$ 是与 $t$ 的、不在 $e$ 上的那个顶点相关联的直线。则给定非零向量 $v_1\in V_1$，存在唯一非零向量 $v_2\in V_2$，使得当 $t$ 位于 $e$ 的右侧时 $v_1+v_2\in V_3$，或当 $t$ 位于 $e$ 的左侧时 $v_2-v_1\in V_3$，见图 10.6。对一条蛇的所有边迭代此过程，便可以对一条蛇赋予一族 $n$ 个向量，它们只依赖于第一个向量的选取；因此这个族在把所有向量同乘一个标量之下是唯一的。

\Fig[0.75\linewidth]{figures/fig_10_6.png}{10.6}{向量 $v_2$ 由 $v_1$ 与三条直线 $(V_1,V_2,V_3)$ 唯一确定。}

\begin{proposition}[10.26] 设 $(F^A,F^B,F^C)\in\mathrm{Conf}^3(\mathcal{F}(\mathbb{K}^n))$ 是横截旗三元组。设 $s$ 是 $S_n$ 上的一条蛇，$(v_1,\dots,v_n)$ 是上述构造的向量族。则 $(v_1,\dots,v_n)$ 是一个射影基。
\end{proposition}

\begin{proof}
不妨设 $s$ 以 $A$ 为基点。则 $v_1$ 是 $F^A_1$ 的一组基；而若 $(v_1,\dots,v_k)$ 是 $F^A_k$ 的基，则由于 $v_{k+1}\in F^A_{k+1}\setminus F^A_k$，$(v_1,\dots,v_{k+1})$ 是 $F^A_{k+1}$ 的基。归纳即得结论。
\end{proof}

% [原文 p.105]
\begin{example}[10.27] 设 $s$ 是图 10.7 中所画的 $S_4$ 上以 $A$ 为基点的蛇。设 $v_1\in A_1$ 非零。向量 $v_2\in A_2\cap B_3$ 由关系 $v_2-v_1\in A_2\cap C_3$ 唯一确定，向量 $v_3\in A_3\cap B_3\cap C_3$ 然后由 $v_2+v_3\in A_3\cap B_2$ 确定，最后 $v_4\in B_2\cap C_3$ 由 $v_4-v_3\in B_3\cap C_2$ 确定。
\end{example}

\Fig[0.5\linewidth]{figures/fig_10_7.png}{10.7}{与一条蛇相伴的射影基的一个例子。}

现在描述两类基本的蛇移动，以及它们所诱导的相伴射影基的变换。第一类移动我们称为 I 型，画在图 10.8 中。这种移动只能在蛇的末端进行，因为它改变终点。第二类移动称为 II 型，也画在图 10.8 中。

\Fig[0.9\linewidth]{figures/fig_10_8.png}{10.8}{左：I 型移动。右：II 型移动。}

容易看出，以 $S_n$ 同一顶点为基点的任意两条蛇都可以通过有限多步 I 型与 II 型移动相联系；关于从 $s_{AB}$ 到 $s_{AC}$ 的移动序列的例子，见定理 10.28 的证明。
现在我们描述这些变换对相伴射影基的影响。考虑图 10.8 左图所画的 $S_2$ 上的蛇。与 $s_{AB}$ 相伴的射影基 $(v_1,v_2)$ 由 $v_1\in A_1$ 的选取以及关系 $v_1+v_2\in C_1$ 确定。于是与 $s_{AC}$ 相伴的射影基为 $(v_1,v_1+v_2)$，这一移动所对应的矩阵是
\[
M_I=\begin{pmatrix}1&1\\0&1\end{pmatrix}.
\]

% [Tutor Note: 这里的矩阵按「基向量排成一行、右乘矩阵」的约定作用：$(v_1,v_2)\mapsto(v_1,v_1+v_2)$ 对应 $(v_1,v_2)\begin{pmatrix}1&1\\0&1\end{pmatrix}=(v_1,\ v_1+v_2)$。核对方法：把两个基向量按行排开相乘即可；若误写成左乘，形状为 $2\times2$ 时结果就会变成 $(v_1+v_2,v_2)$，与图中把终点从 $B$ 换到 $C$ 的情形不符。]

% [原文 p.106]
现在考虑图 10.8 右图中的 II 型移动。设 $(v_1,v_2,v_3)$ 是与第一条蛇相伴的射影基。则与第二条蛇相伴的射影基为 $(v_1,v'_2,v'_3)$，其中 $v'_2=v_1+v_2$；又由于 $v_3$ 与 $v'_3$ 属于同一条直线，故 $v'_3=\lambda v_3$，其中 $\lambda\ne0$。事实上，$\lambda$ 正是该旗三元组的三重比的倒数。确实，取如下向量：$v_1\in A_1$、$v_2\in A_2/A_1$、$b_1=v_3-v_2\in B_1$、$v_3\in B_2/B_1$、$c_1=v'_2+v'_3\in C_1$ 以及 $v'_2\in C_2/C_1$。则我们有
\begin{align*}
r(A,B,C) &= \frac{\omega(v_1,v_2,b_1)\omega(b_1,v_3,c_1)\omega(c_1,v'_2,v_1)}{\omega(v_1,v_2,c_1)\omega(b_1,v_3,v_1)\omega(c_1,v'_2,b_1)}\\
&= \frac{\omega(v_1,v_2,-v_3)\omega(v_2,v_3,c_1)\omega(c_1,v_2,v_1)}{\omega(v_1,v_2,c_1)\omega(v_2,v_3,v_1)\omega(v'_3,v'_2,b_1)}\\
&= \frac{\omega(v_2,v_3,c_1)\omega(c_1,v_2,v_1)}{\omega(v_1,v_2,c_1)\omega(v'_3,v'_2,b_1)}\\
&= \frac{-\omega(v_2,v_3,c_1)}{\omega(v'_3,v'_2,b_1)}\\
&= \frac{-\omega(v_2,v_3,c_1)}{\omega(v'_3,v'_2,v_2)}\\
&= \frac{-\omega(v_2,v_3,c_1)}{\omega(v'_3,c_1,v_2)}\\
&= \frac{1}{\lambda}.
\end{align*}

% [Tutor Note: 原文此处由第四行到第五行把 $\omega(v'_3,v'_2,b_1)$ 换成了 $\omega(v'_3,v'_2,v_2)$，但由 $b_1=v_3-v_2$ 与 $v'_3=\lambda v_3$ 有 $\omega(v'_3,v'_2,b_1)=\omega(v'_3,v'_2,v_3)-\omega(v'_3,v'_2,v_2)=-\omega(v'_3,v'_2,v_2)$，故正确的替换会使分数变号；按正确符号一路算到末行恰得 $1/\lambda$，与正文「$\lambda$ 是三重比的倒数」一致。这里的印刷符号有一处遗漏，照录原文，计算时以结论 $r(A,B,C)=1/\lambda$ 为准。]

于是这一变换对应的矩阵是
\begin{align*}
M_{II}&=\begin{pmatrix}1&1&0\\0&1&0\\0&0&r(A,B,C)^{-1}\end{pmatrix}
=\begin{pmatrix}1&1&0\\0&1&0\\0&0&1\end{pmatrix}
\begin{pmatrix}1&0&0\\0&1&0\\0&0&r(A,B,C)^{-1}\end{pmatrix}\\
&=\begin{pmatrix}1&0&0\\0&1&0\\0&0&r(A,B,C)^{-1}\end{pmatrix}
\begin{pmatrix}1&1&0\\0&1&0\\0&0&1\end{pmatrix}.
\end{align*}
现在我们处理一般情形：对 $i\le n$，设 $\varphi_i:\SL_2(\mathbb{K})\to\GL_n(\mathbb{K})$ 是对应于第 $i$ 个单根的同态。令
\[
E_i:=\varphi_i\begin{pmatrix}1&1\\0&1\end{pmatrix},
\qquad
H_i(x):=\diag(\underbrace{x,\dots,x}_{i},1,\dots,1).
\]

% [原文 p.107]
注意当 $i\ne j$ 时 $H_i(x)$ 与 $E_j$ 交换。

\begin{theorem}[10.28] 设 $(A,B,C)$ 是 $\mathbb{K}^n$ 中的横截旗三元组，且对满足 $i+j+k=2n+3$ 的 $i,j,k\ge1$，设 $r_{i,j,k}$ 为 $(A,B,C)$ 在 $A_i\cap B_j\cap C_k$ 上诱导的旗三元组的三重比。设 $M_{AB\to AC}$ 是与把 $s_{AB}$ 变换为 $s_{AC}$ 相对应的矩阵，$M_{AB\to BA}$ 是与把 $s_{AB}$ 变换为 $s_{BA}$ 相对应的矩阵。则
\[
M_{AB\to AC}=\left[\prod_{j=n-1}^{1}\left[\prod_{i=j}^{n-2}H_{i+1}(r_{n-i+1,\,n-j+1,\,i+j+1})E_i\right]E_{n-1}\right]
\]
而 $M_{AB\to BA}$ 由 $M_{AB\to BA}(e_i)=(-1)^{i-1}e_{n+1-i}$ 给出。
\end{theorem}

\begin{proof}
对第一点，最内层的乘积来自图 10.9 所描述的移动序列，最外层的乘积对应于图 10.10 的移动序列。
对第二点，设 $(v_1,\dots,v_n)$ 是与 $s_{AB}$ 相伴的射影基，$(v'_1,\dots,v'_n)$ 是与 $s_{BA}$ 相伴的射影基。适当地乘以标量，可设 $v'_1=v_n$。于是我们有 $v'_2-v'_1=v_{n-1}+v_n$，从而 $v_{n-1}=-v'_2$。归纳即得结论。
\end{proof}

% [Tutor Note: 定理 10.28 中 $M_{AB\to BA}$ 可显式写出：由 $M(e_i)=(-1)^{i-1}e_{n+1-i}$ 得 $M_{AB\to BA}=\begin{pmatrix}0&\cdots&0&1\\ \vdots&&1&0\\ 0&\ddots&&\vdots\\ -1&0&\cdots&0\end{pmatrix}$，即第 $i$ 列在第 $n+1-i$ 行处为 $(-1)^{i-1}$。核对 $n=2$：$M_{AB\to BA}=\begin{pmatrix}0&1\\-1&0\end{pmatrix}$，它正好把 $(v_1,v_2)$ 送到 $(v_2,-v_1)$。另外，证明第二点里由 $v'_2-v'_1=v_{n-1}+v_n$ 与 $v'_1=v_n$ 并不能直接推出 $v_{n-1}=-v'_2$（代入得 $v'_2=v_{n-1}+2v_n$），此处显然是原文的一处符号/下标笔误；正确的归纳结论应为 $v'_i=(-1)^{i-1}v_{n+1-i}$，而 $v'_2=-v_{n-1}$。]

\Fig[0.8\linewidth]{figures/fig_10_9.png}{10.9}{诱导出定理 10.28 中最内层乘积所给变换的移动序列。}

设 $(V_1,V_2,V_3,V_4)$ 是 $\mathbb{K}^2$ 中的四条横截直线。设 $s$ 是三角形 $(V_1,V_2,V_3)$ 中 $S_2$ 上从 $V_1$ 到 $V_2$ 的蛇，$s'$ 是三角形 $(V_1,V_4,V_2)$ 中 $S_2$ 上从 $V_1$ 到 $V_2$ 的蛇。设 $(v_1,v_2)$ 是与 $s$ 相伴的射影基，$(v_1,v'_2)$ 是与 $s'$ 相伴的射影基，其中 $v'_2=\lambda v_2$。按定义，
% [原文 p.108]
\Fig[0.8\linewidth]{figures/fig_10_10.png}{10.10}{诱导出定理 10.28 中最外层乘积所给变换的移动序列。}
$v_1+v_2\in V_3$ 且 $v'_2-v_1\in V_4$。设 $\omega$ 是 $\mathbb{K}^2$ 上的体积形式。该四元组直线的交比为
\begin{align*}
x &= \frac{\omega(v_1,v_1+v_2)\omega(v_2,v'_2-v_1)}{\omega(v_1,v'_2-v_1)\omega(v_1+v_2,v_2)}\\
&= \frac{\omega(v_1,v_2)\omega(v_2,-v_1)}{\omega(v_1,v'_2)\omega(v_1,v_2)}\\
&= \frac{1}{\lambda}.
\end{align*}

\Fig[0.4\linewidth]{figures/fig_10_11.png}{10.11}{一条蛇的一条边属于两个不同的三角形。}

对一条蛇的每条属于两个不同三角形的边重复这一论证（图 10.11），便得到下述结果：

\begin{theorem}[10.29] 设 $(A,B,C)$ 与 $(A,D,B)$ 是两个横截旗三元组。设 $s_{T^+(AB)}$ 是三角形 $(A,B,C)$ 中从 $A$ 到 $B$ 的蛇，$s_{T^-(AB)}$ 是三角形 $(A,D,B)$ 中从 $A$ 到 $B$ 的蛇。设 $(v_1,\dots,v_n)$ 是与 $s_{T^+(AB)}$ 相伴的射影基，$(v'_1,\dots,v'_n)$ 是与
% [原文 p.109]
$s_{T^-(AB)}$ 相伴的射影基。设 $x_1,\dots,x_{n-1}$ 是依次围绕 $s_{AB}$ 的四元组直线的交比。则我们有
\[
M_{T^+(AB)\to T^-(AB)}=\diag(1,x_1^{-1},\dots,x_1^{-1}\cdots x_{n-1}^{-1}).
\]
\end{theorem}

现在我们可以由与 $S$ 的三角剖分 $\Delta$ 相关联的 X-坐标的数据，重构一个带标架表示 $\rho:\pi_1(S)\to\PGL_n(\mathbb{K})$。注意所有与蛇的变换相关联的矩阵都只依赖于与所选三角剖分相关联的 X-坐标。我们首先通过在 $\Delta$ 的每个三角形上粘合基本构件来构造一个图 $\Gamma$，见图 10.12。我们给 $\Gamma$ 的每条有向边赋一个可逆矩阵，它对应于如图 10.12 所示的蛇的变换。现在 $S$ 上的任何环路 $\gamma$ 都可以收缩到图 $\Gamma$ 上，它在 $\rho$ 下的像将是 $\gamma$ 所经过的 $\Gamma$ 的边对应矩阵的乘积。所得的矩阵不依赖于把 $\gamma$ 收缩到 $\Gamma$ 的方式，因为绕 $\Gamma$ 的六边形面与四边形面的和乐都是平凡的。与从某个穿孔出发的任何蛇相伴的射影基诱导出一个旗，因此所构造的表示自然地带有标架。

\Fig[0.75\linewidth]{figures/fig_10_12.png}{10.12}{用来在 $S$ 上重构带标架局部系统的图 $\Gamma$ 上的局部系统。}

% [Tutor Note: 重构的良定义性可以这样核对：$\Gamma$ 的每条边上的矩阵都是射影基之间的过渡矩阵，因此整体表示在把某个顶点处所有基向量同乘一个标量下不变；绕六边形面的和乐平凡正是定理 10.28 中 $M_{AB\to AC}$ 与其逆矩阵配对的体现，绕四边形面平凡则对应定理 10.29 的对角矩阵 $M_{T^+(AB)\to T^-(AB)}$ 与其逆配对（对角元恰为各交比的倒数）。可用 $n=2$ 的图 10.12 直接检验：此时每个三角形里的矩阵只有 $M_I$ 型与 $M_{II}$ 型，绕一个六边形走一圈的乘积为 $I$。]

% [原文 p.110]
% 原书此页为空白页（第 10 章正文结束于印刷 p.109 / PDF p.149）。
